\chapter{Point Location}
\label{ch:point_location}

% Stub: outline only; content to be written.

\section{The Planar Point-Location Problem}
\label{sec:pointloc_problem}

\section{Subdivision Representations}
\label{sec:pointloc_subdivisions}

\section{Slab Decomposition}
\label{sec:pointloc_slabs}

\section{Trapezoidal Maps}
\label{sec:pointloc_trapezoidal}

\section{Randomized Incremental Construction}
\label{sec:pointloc_randomized}

\section{Search DAGs}
\label{sec:pointloc_search_dags}

\section{Kirkpatrick's Hierarchy}
\label{sec:pointloc_kirkpatrick}

\section{Point Location in Triangulations}
\label{sec:pointloc_triangulations}

Walking methods locate a query point by crossing mesh edges (2D) or faces
(3D) from a starting simplex toward the query point. The visibility walk and
jump-and-walk are described in
Section~\ref{sec:point_location_walking}, which covers both triangles and
tetrahedra and the expected $O(\log n)$ walk length with a spatial-index
start.

\section{Dynamic Point Location}
\label{sec:pointloc_dynamic}

\section{Practical Implementations}
\label{sec:pointloc_practical}

\section{Exercises}
\label{sec:pointloc_exercises}
